\documentclass[12pt, openright, twoside, a4paper, english, french, spanish]{abntex2}

% ---
% Pacotes básicos
% ---
\usepackage{lmodern}            % Usa a fonte Latin Modern
\usepackage[T1]{fontenc}        % Selecao de codigos de fonte.
\usepackage[utf8]{inputenc}     % Codificacao do documento (conversão automática dos acentos)
\usepackage{indentfirst}        % Indenta o primeiro parágrafo de cada seção.
\usepackage{color}              % Controle das cores
\usepackage{graphicx}           % Inclusão de gráficos
\usepackage{microtype}          % para melhorias de justificação
% ---

% ---
% Pacotes de citações
% ---
\usepackage[brazilian,hyperpageref]{backref}     % Paginas com as citações na bibl
\usepackage[alf]{abntex2cite}   % Citações padrão ABNT

% ---
% Informações de dados para CAPA e FOLHA DE ROSTO
% ---
\titulo{Soberania de Dados e Auditoria Algorítmica em Gestão Clínica: O Caso Intraclinica}
\autor{Bruno Banho \and Axiomatic Delirium (Co-Autor Digital)}
\local{São Paulo, Brasil}
\data{2026}
\instituicao{%
  Axio Engineering
  \par
  Divisão de Sistemas Soberanos
}
\tipotrabalho{Artigo Técnico}
% ---

% ---
% Configurações de aparência do PDF final
% ---
\definecolor{blue}{RGB}{41,5,195}
\makeatletter
\hypersetup{
    pdftitle={\@title}, 
    pdfauthor={\@author},
    pdfsubject={\@title},
    pdfcreator={LaTeX with abntex2},
    pdfkeywords={Gestão Clínica}{Auditoria}{IA}{Soberania de Dados}, 
    colorlinks=true,       % false: box links; true: colored links
    linkcolor=blue,        % color of internal links
    citecolor=blue,        % color of links to bibliography
    filecolor=magenta,     % color of file links
    urlcolor=blue,
    bookmarksdepth=4
}
\makeatother
% ---

% ---
% Espaçamentos entre linhas e parágrafos 
% ---
\setlength{\parindent}{1.3cm} % O tamanho do parágrafo é dado por 1.3cm
\setlength{\parskip}{0.2cm}  % Controle do espaçamento entre um parágrafo e outro

% ---
% Início do documento
% ---
\begin{document}

% Seleciona o idioma principal do documento (para hifenização e palavras-chave)
\selectlanguage{brazil}

% Retira espaço extra obsoleto entre as frases.
\frenchspacing 

% ----------------------------------------------------------
% ELEMENTOS TEXTUAIS
% ----------------------------------------------------------
\textual

% ----------------------------------------------------------
% Introdução
% ----------------------------------------------------------
\section*{Introdução}
\addcontentsline{toc}{section}{Introdução}

O setor de saúde suplementar no Brasil enfrenta uma crise silenciosa de eficiência operacional. Enquanto a tecnologia médica avança exponencialmente, a gestão administrativa permanece analógica ou refém de sistemas legados monolíticos. Este artigo apresenta a arquitetura e a proposta de valor do sistema \textbf{Intraclinica}, uma solução focada na soberania de dados e na auditoria algorítmica de estoques.

O problema central identificado é a desconexão entre o ato médico e o registro administrativo, gerando o que denominamos de "Margem Misteriosa" --- a perda de lucratividade por insumos consumidos mas não faturados.

% ----------------------------------------------------------
% Desenvolvimento
% ----------------------------------------------------------
\section{O Paradigma da Auditoria Ativa}

Diferente dos sistemas ERP tradicionais, que dependem da imputação manual de dados (passiva), o Intraclinica opera sob o paradigma da Auditoria Ativa. O sistema utiliza inferência lógica baseada em procedimentos para deduzir o consumo de materiais.

\subsection{O Guardião de Estoque (Inventory Guardian)}

A arquitetura do sistema vincula procedimentos médicos a receitas de consumo estritas. Ao registrar um procedimento clínico (e.g., Sutura Simples), o sistema dispara um gatilho de banco de dados (RPC) que executa a baixa automática dos insumos correlatos (Fio de Nylon, Kit de Pequena Cirurgia, Luvas Estéreis).

Isso elimina a necessidade de a equipe de enfermagem realizar a baixa manual item a item, reduzindo a carga cognitiva e o erro humano.

\section{Soberania de Dados e Arquitetura}

Em um cenário onde dados de saúde são frequentemente monetizados por grandes plataformas, o Intraclinica adota uma postura de \textit{Tenant Isolation} rigorosa. Utilizando \textit{Row Level Security} (RLS) no nível do banco de dados (PostgreSQL/Supabase), garantimos que os dados de cada clínica sejam criptograficamente isolados, mesmo em uma infraestrutura multi-tenant.

\section{Metodologia de Validação e Métricas Esperadas}

O sistema encontra-se em fase de validação pré-operacional. O protocolo de teste piloto, desenhado para uma clínica parceira de médio porte, visa mensurar dois indicadores chave de desempenho (KPIs):

\begin{enumerate}
    \item \textbf{Delta de Inventário (Divergência):} A diferença percentual entre o estoque físico e o estoque lógico deduzido pelo algoritmo após 30 dias de operação.
    \item \textbf{Tempo de Ocupação Administrativa:} A medição cronometrada do tempo gasto pela equipe de enfermagem no lançamento de insumos, comparando o método atual (manual) versus o método proposto (automático por procedimento).
\end{enumerate}

A hipótese central deste trabalho é que a automação da baixa de estoque reduzirá a divergência de inventário para patamares inferiores a 5\%, transformando a gestão de insumos de um centro de custo passivo para um ativo auditável em tempo real.

% ----------------------------------------------------------
% Conclusão
% ----------------------------------------------------------
\section*{Considerações Finais}
\addcontentsline{toc}{section}{Considerações Finais}

A arquitetura Intraclinica propõe uma mudança fundamental na forma como softwares médicos interagem com a realidade física do consultório. Ao priorizar a soberania de dados e a automação lógica, o sistema busca devolver ao médico o controle sobre sua margem operacional. Os próximos passos envolvem a coleta de dados empíricos para validar a eficácia do modelo de "Guardião de Estoque" em cenário real.

\end{document}
